% answers_short.tex —— 2015 年·浙江卷·理科 · 简短答案速查
% 由 tools/make_short_answers.py 自 answers.tex 的 \gksolhead/\gkanswer 提取，
% 并与 2015/2015_数学真题_浙江理_答案要点.md 逐题比对后生成，勿手改（改源 MD 后重跑）。

\begin{gknotice}
\textbf{使用说明}：本节只给最终答案便于快速核对，完整推导见解析版。答案为编写者独立推导，\textbf{非官方标准答案}。
\end{gknotice}

\gkpart{选择题}
\begin{center}\begin{tabular}{ccccccccc}
\toprule
题号 & 1 & 2 & 3 & 4 & 5 & 6 & 7 & 8 \\
\midrule
答案 & C & C & B & D & A & A & D & B \\
\bottomrule
\end{tabular}\end{center}

\gkpart{填空题}
\begin{itemize}[leftmargin=2.6em]
  \item[\textbf{9.}] \(2\sqrt3\)；\(y=\pm\dfrac{\sqrt2}{2}x\)
  \item[\textbf{10.}] \(0\)；\(2\sqrt2-3\)
  \item[\textbf{11.}] \(\pi\)；\(\left[\dfrac{3\pi}{8}+k\pi,\dfrac{7\pi}{8}+k\pi\right]\)（\(k\in\mathbf{Z}\)）
  \item[\textbf{12.}] \(\dfrac{4\sqrt3}{3}\)
  \item[\textbf{13.}] \(\dfrac78\)
  \item[\textbf{14.}] \(3\)
  \item[\textbf{15.}] \(x_0=1\)，\(y_0=2\)，\(|\boldsymbol{b}|=2\sqrt2\)
\end{itemize}

\gkpart{解答题}
\begin{itemize}[leftmargin=2.6em]
  \item[\textbf{16.}] (1) \(2\)；(2) \(3\)
  \item[\textbf{17.}] (1) 成立（建系得 \(\overrightarrow{A_1D}\) 与平面 \(A_1BC\) 内两条相交直线 \(A_1B\)、\(A_1C\) 均垂直）；(2) \(-\dfrac18\)
  \item[\textbf{18.}] (1) 成立（\(M(a,b)\geqslant\max\{|f(1)|,|f(-1)|\}\geqslant\dfrac12|f(1)-f(-1)|=|a|\geqslant2\)）；(2) \(3\)
  \item[\textbf{19.}] (1) \(m\in\left(-\infty,-\dfrac{\sqrt6}{3}\right)\cup\left(\dfrac{\sqrt6}{3},+\infty\right)\)；(2) \(\dfrac{\sqrt2}{2}\)
  \item[\textbf{20.}] (1) 成立（\(\dfrac{a_n}{a_{n+1}}=\dfrac{1}{1-a_n}\)，由 \(0<a_n\leqslant\dfrac12\) 得 \(\dfrac{1}{1-a_n}\in[1,2)\)）；(2) 成立（\(S_n=\dfrac12-a_{n+1}\)，结合 \(\dfrac{1}{2k}\leqslant a_k\leqslant\dfrac{1}{k+1}\)（\(k\geqslant2\)）即得）
\end{itemize}
